% Part: first-order-logic
% Chapter: natural-deduction
% Section: provability-propositional

% verification of properties of provability needed for maximally
% consistent sets in the completeness chapter.

\documentclass[../../../include/open-logic-section]{subfiles}

\begin{document}

\iftag{FOL}
      {\olfileid{fol}{ntd}{ppr}}
      {\olfileid{pl}{ntd}{ppr}}

\olsection{\usetoken{S}{derivability} మరియు ప్రతిపాదనాత్మక సంయోజకాలు}

\begin{explain}
  సహజ నిగమనపు !!{derivability} సంబంధం~$\Proves$, $!A \land !B
  \Proves !A$, $!A, !A \lif !B \Proves !B$ (మోడస్ పోనెన్స్) వంటి
  ప్రతిపాదనాత్మక సంయోజకాల ప్రాథమిక వాస్తవాలను స్థాపించగలిగేంత
  బలమైనదని చూపుతాం. సంపూర్ణతా సిద్ధాంతపు నిరూపణకు ఈ వాస్తవాలు అవసరం.
\end{explain}

\begin{prop}\ollabel{prop:provability-land}
  \begin{enumerate}
  \item \ollabel{prop:provability-land-left} $!A \land !B \Proves
    !A$, $!A \land !B \Proves !B$ రెండూ.
  \item \ollabel{prop:provability-land-right} $!A, !B \Proves !A \land !B$.
  \end{enumerate}
\end{prop}

\begin{proof}
  \begin{enumerate}
  \item రెండింటినీ ఇలా !!{derive} చేయవచ్చు:
    \begin{prooftree}
      \AxiomC{$!A \land !B$}
      \RightLabel{\Elim{\land}}
      \UnaryInfC{$!A$}
      \DisplayProof\qquad\bottomAlignProof
      \AxiomC{$!A \land !B$}
      \RightLabel{\Elim{\land}}
      \UnaryInfC{$!B$}
    \end{prooftree}
  \item ఇలా !!{derive} చేయవచ్చు:
    \begin{prooftree}
      \AxiomC{$!A$}
      \AxiomC{$!B$}
      \RightLabel{\Intro{\land}}
      \BinaryInfC{$!A \land !B$}
    \end{prooftree}
  \end{enumerate}
\end{proof}


\begin{prop}\ollabel{prop:provability-lor}
  \begin{enumerate}
  \item $!A \lor !B, \lnot !A, \lnot !B$ వైరుధ్యభరితమైనది.
  \item $!A \Proves !A \lor !B$, $!B \Proves !A \lor !B$ రెండూ.
  \end{enumerate}
\end{prop}

\begin{proof}
  \begin{enumerate}
  \item కింది !!{derivation}ను పరిశీలించండి:
    \begin{prooftree}
      \AxiomC{$!A \lor !B$}
      \AxiomC{$\lnot !A$}
      \AxiomC{$\Discharge{!A}{1}$}
      \RightLabel{\Elim{\lnot}}
      \BinaryInfC{$\lfalse$}
      \AxiomC{$\lnot !B$}
      \AxiomC{$\Discharge{!B}{1}$}
      \RightLabel{\Elim{\lnot}}
      \BinaryInfC{$\lfalse$}
      \DischargeRule{\Elim{\lor}}{1}
      \TrinaryInfC{$\lfalse$}
    \end{prooftree}
    ఇది !!{undischarged} పరికల్పనలు $!A \lor !B$, $\lnot !A$,
    $\lnot !B$ల నుంచి $\lfalse$ యొక్క ఒక !!a{derivation}.
  \item రెండింటినీ ఇలా !!{derive} చేయవచ్చు:
    \begin{prooftree}
      \AxiomC{$!A$}
      \RightLabel{\Intro{\lor}}
      \UnaryInfC{$!A \lor !B$}
      \DisplayProof\qquad\bottomAlignProof
      \AxiomC{$!B$}
      \RightLabel{\Intro{\lor}}
      \UnaryInfC{$!A \lor !B$}
    \end{prooftree}
  \end{enumerate}
\end{proof}

\begin{prop}\ollabel{prop:provability-lif}
  \begin{enumerate}
  \item \ollabel{prop:provability-lif-left} $!A, !A \lif !B \Proves !B$.
  \item \ollabel{prop:provability-lif-right}
    $\lnot !A \Proves !A \lif !B$, $!B \Proves !A \lif !B$ రెండూ.
  \end{enumerate}
\end{prop}

\begin{proof}
  \begin{enumerate}
  \item ఇలా !!{derive} చేయవచ్చు:
    \begin{prooftree}
      \AxiomC{$!A \lif !B$}
      \AxiomC{$!A$}
      \RightLabel{\Elim{\lif}}
      \BinaryInfC{$!B$}
    \end{prooftree}

  \item కింది రెండు !!{derivation}s దీన్ని చూపుతాయి:
    \begin{prooftree}
      \AxiomC{$\lnot !A$}
      \AxiomC{$\Discharge{!A}{1}$}
      \RightLabel{\Elim{\lnot}}
      \BinaryInfC{$\lfalse$}
      \RightLabel{\FalseInt}
      \UnaryInfC{$!B$}
      \DischargeRule{\Intro{\lif}}{1}
      \UnaryInfC{$!A \lif !B$}
      \DisplayProof\qquad\bottomAlignProof
      \AxiomC{$!B$}
      \RightLabel{\Intro{\lif}}
      \UnaryInfC{$!A \lif !B$}
    \end{prooftree}
    $\Intro{\lif}$ పరికల్పన~$!A$ను !!{discharge} చేయవచ్చు; కానీ
    తప్పనిసరిగా చేయనవసరం లేదని గమనించండి.
  \end{enumerate}
\end{proof}

\end{document}
